\documentclass[instructions.tex]{subfiles}
\begin{document}
 
\chapter{De la patience dans les maladies.}
 
Souffrez la maladie avec résignation, parce que Dieu le veut; avec consolation, parce qu’elle est avantageuse.
 
\primo Dieu le veut. Pourquoi ne le voudriez-vous pas ? Qu’il y a de bizarrerie dans notre esprit ! Tous les jours nous disons à Dieu : Votre volonté soit faite. Mais, s’il manifestait sa volonté par quelque disgrâce, nous voudrions que sa volonté ne se fît plus. Dieu pourrait vous préserver de maladie ou vous guérir. S’il ne le fait pas, c’est une marque qu’il veut que vous souffriez; vous devez donc le vouloir. Soulagez-vous par des remèdes, il le permet, il le veut; mais remettez à sa volonté l’opération et l’effet des remèdes. En vain résistez-vous à sa volonté toute-puissante; faites du moins de nécessité vertu.
 
Vous n’êtes pas le seul que Dieu afflige. Combien d’autres sont accablés de peines ! Y a-t-il sur la terre un homme sans disgrâce, une famille sans adversité ? N’a-t-on pas ses chagrins sur le trône aussi bien que dans les fers ? Si tant d’autres sont obligés à porter leur croix, pourquoi n’acceptez-vous pas la vôtre ?
 
Vous vous croyez plus affligé que les autres; mais ne vous trompez-vous point ? ils sont peut-être plus malheureux que vous. Vos amis pensent que vous êtes à plaindre; vous ne l’êtes pas parce qu’on le pense, mais vous l’êtes si vous souffrez sans résignation. Vos souffrances fussent-elles encore plus accablantes, jamais la maladie, ni rien au monde, ne doit vous empêcher d’être soumis à Dieu. Quand vous manqueriez de toute consolation, vous devez être content de ce que Dieu veut.
 
Cette divine volonté est si adorable et si sainte, que, quand il ne faudrait qu’un pas pour vous délivrer, vous ne devriez pas le faire. Il est plus sûr pour vous d’être malade, selon la volonté de Dieu, que d’être en santé contre son gré. Tels étaient les sentimens d’une vertueuse villageoise, qui, souffrant une cruelle maladie, me dit un jour : Je suis si contente d’être ce que Dieu veut, que je ne changerois pas mes infirmités contre un royaume. S’inquiéter parce que la maladie empêche de remplir ses devoirs et de servir Dieu, c’est une illusion. Pourquoi vous inquiéter de ne pouvoir faire ce que Dieu ne demande pas de vous ? Qu’y a-t-il de plus grand, et qui donne plus de gloire à Dieu, que de souffrir pour son amour ?
 
\secundo Il y a d’ailleurs de grands sujets de consolation dans une maladie, par les avantages qui l’accompagnent. Elle est avantageuse, 1. pour ceux qui vous assistent; ils gagnent des mérites immenses par leur charité. Oh ! que je suis consolé, disoit saint François de Sales dans ses maladies, de voir la peine que ces pauvres gens prennent autour de moi ! Par leurs services et leur charité ils gagnent le ciel.
 
2. Avantageuse pour l’expiation de vos péchés. Il est bien plus doux et plus court d’effacer ses péchés sur un lit que de les expier dans le feu. Combien d’années de purgatoire ne peut-on pas expier dans quelques heures de maladie ! Ah ! qu’un jour vous vous saurez bon gré de vous voir délivré de ces feux horribles, dont une heure est plus insupportable que plusieurs années de la maladie la plus aiguë !
 
3. Avantageuse pour votre prédestination. Il est peu de moyens plus efficaces pour purifier une âme et pour la sauver, qu’une longue maladie. Plusieurs grands saints ont passé toute leur vie dans des infirmités habituelles. Ils ne pouvaient vaquer à l’oraison, ni à plusieurs actes de religion; ils étaient soumis à Dieu; ils souffraient pour son amour : voilà tout ce que Dieu demandait pour les sanctifier.
 
Saint Alype, solitaire, demeura couché sur un côté tout écorché, pendant quatorze ans. Dans cette cruelle situation, sa prière était : J’adore votre volonté sainte, ô mon Dieu ! vous êtes juste, et vous me punissez avec justice. Sainte Ludivine pria le Seigneur de lui ôter sa beauté; elle fut exaucée. Dieu l’affligea par une maladie extraordinaire pendant trente-huit ans. Sa patience l’éleva à une haute sainteté. Oh ! combien sont en enfer, qui seraient dans le ciel, s’ils avaient été long-temps malades et affligés !
 
Pourvu que vous alliez au ciel, il importe peu par quel chemin. Si ce chemin vous paraît long, il vous paraîtra bien court quand vous serez au terme. Laissez agir la Providence, et mettez en Dieu votre confiance. Ceux qui passent la mer se fient à leur pilote pendant la tempête, sans se plaindre de sa conduite. Cette vie est une mer orageuse qu’il faut passer parmi les écueils. Auriez-vous moins de confiance en Dieu pour vous conduire au ciel, qu’on en a en un pilote pour arriver au port ?
 
Dieu sait ce qui vous convient pour votre salut. Contentez-vous de le laisser agir et de le suivre. Il n’a pas besoin de conseil; il ne demande que votre soumission. Sa bonté a disposé vos croix selon vos forces; il a fixé le temps de vous en délivrer : elles ne dureront qu’autant qu’il est nécessaire pour votre salut. Consolez-vous dans cette pensée, que rien ne vous arrive que Dieu n’ait ordonné pour votre prédestination.
 
Après tout, Dieu est le souverain maître; les rois, les monarques, tous les hommes, sont devant lui comme un grain de poussière. Nous sommes son ouvrage, il peut nous conserver ou nous réduire en poudre, sans rendre compte de sa conduite, qui est toujours adorable et pleine d’équité. Je recevrai donc, ô mon Dieu ! dit l’auteur de l’Imitation, les biens et les maux de votre main paternelle. Préservez-moi du péché; pourvu que vous ne me rejetiez pas dans l’éternité, que vous ne m’effaciez pas du livre de vie, quelque tribulation qui m’arrive, rien ne pourra me nuire (1).\footnote{liv. 3. 17.}
 
\section{Résolutions}
 
\primo Je recevrai les maladies qui m’arriveront, comme me venant de la main de Dieu.
 
\secundo Et je m’y soumettrai en esprit de pénitence, et pour reconnoître ma dépendance envers le souverain arbitre de la vie et de la mort. 


Mon Dieu, je m’abandonne à votre sainte volonté pour la santé et la maladie; faites que je vous demeure fidèle dans toutes les situations par lesquelles il vous plaira de me faire passer.
 
\end{document}
